\section{Renormalization of lattice smeared matrix elements}\label{sec:smearing}
  
In many lattice calculations, the data can be very noisy
% due to large fluctuations,
which calls for more statistics. However, 
large statistics means more resources which are hard to 
come by sometimes. Therefore, lattice practitioners have
invented phenomenological approaches to quench the short-distance
fluctuations: smearing. However, smearing might lead 
to configurations that are not connected with fundamental theory. On the other hand, smearing in some sense just makes the effective UV cutoff (the inverse of the lattice spacing) smaller, and in the limit $a\to 0$, all smearing becomes local effects which can be taken into account properly through some kind of renormalization. 
In this section, we check if the procedure we discussed
in the previous section still works for smeared matrix elements. 



\begin{figure*}[tbp]
\centering
\includegraphics[width=8.8cm]{images/fitting_overlap_quark_hyp1_Oa.pdf}
\includegraphics[width=8.8cm]{images/fitting_clover_quark_hyp1_Oa.pdf}
\includegraphics[width=8.8cm]{images/fitting_overlap_Pion_hyp1_Oa.pdf}
\includegraphics[width=8.8cm]{images/fitting_clover_Pion_hyp1_Oa.pdf}
\includegraphics[width=8.8cm]{images/fitting_clover_Nucleon_hyp1_Oa.pdf}
\caption{Using Eq.~(\ref{eq:logM_ss}) to fit $O_{\gamma_t}(z)$ matrix element for HYP smearing cases. $\Lambda_{\rm QCD}$ is fixed at 0.39 GeV.
}
\label{fig:fitting_hyp}
\end{figure*}


\begin{figure*}[tbp]
\centering
\includegraphics[width=8.5cm]{images/ez_vs_z_overlap_quark_hyp1_Oa.pdf}
\includegraphics[width=8.5cm]{images/ez_vs_z_clover_quark_hyp1_Oa.pdf}
\includegraphics[width=8.5cm]{images/ez_vs_z_overlap_Pion_hyp1_Oa.pdf}
\includegraphics[width=8.5cm]{images/ez_vs_z_clover_Pion_hyp1_Oa.pdf}
\includegraphics[width=8.5cm]{images/ez_vs_z_clover_Nucleon_hyp1_Oa.pdf}
\caption{$e(z)$ with respect to $z$. $e(z)$ is extracted from the fitting in Fig.~\ref{fig:fitting_hyp}. The slopes are labeled in the plot legend.
}
\label{fig:e(z)_hyp}
\end{figure*}

To ensure that the HYP smearing data are useful for refined analysis, we can do a simple test on whether the HYP smearing data show the properties of the linear divergence predicted by perturbation theory, as we have done in Sec.~\ref{sec:test}. Here we use the following function to fit the bare matrix element $\mathcal{M}$ to extract the linear divergence factor,
\begin{align}\label{eq:logM_ss}
\ln \mathcal{M}(z,a) = \frac{e(z)}{a \ln[a \Lambda_{\rm QCD}]} 
+ g(z) + f_{1,2}(z)a,
\end{align}
where we allow for a discretization error to get a better fit. We fix $\Lambda_{\rm QCD}$ at 0.39 GeV in this simple test. Fig.~\ref{fig:fitting_hyp} shows the fits for different actions and states which all look reasonable. 
We have not shown $\chi^2$ which will be considered in the extraction
of the renormalization factor.

Fig.~\ref{fig:e(z)_hyp} shows the extracted $e(z)$ with respect to $z$. The linear $z$ dependence is approximately reproduced but some deviations are seen at large $z$. We can perform a linear fitting of $e(z)$ and compare the fitted slopes.
Overall, the quality of the fits is not as good as in the unsmeared cases, although it is still rather impressive. 
The slope extracted for the linear divergence for the clover quark is totally different from the others, which explains why we cannot use the RI/MOM factor to eliminate it in the clover case for HYP smearing data~\cite{Zhang:2020rsx}. Although the slopes for the overlap quark and pion, clover pion and nucleon are similar to each other, there are some small discrepancies. Using the RI/MOM factor for these cases may leave small residual linear divergences, and the self-renormalization method provides a better way to renormalize the matrix elements. 

Next, we extend our self-renormalization method to the HYP smeared data. Our fitting function, Eq.~(\ref{eq:logM}) may have no clear physical meanings for the HYP smeared date, 
but it may be regarded as a phenomenological model.  

\begin{figure*}[tbp]
\centering
\includegraphics[width=8.5cm]{images/chisqmap_overlap_quark_hyp1.pdf}
\includegraphics[width=8.5cm]{images/chisqmap_clover_quark_hyp1.pdf}
\includegraphics[width=8.5cm]{images/chisqmap_overlap_Pion_hyp1.pdf}
\includegraphics[width=8.5cm]{images/chisqmap_clover_Pion_hyp1.pdf}
\includegraphics[width=8.5cm]{images/chisqmap_clover_Nucleon_hyp1.pdf}
\caption{Using Eq.~(\ref{eq:logM}) to fit the $O_{\gamma_t}(z)$ matrix element for HYP smearing. The plot shows the $\chi^{2}$ map with respect to $k$ and $\Lambda_{\rm QCD}$ for different actions and states. $\langle  \chi^{2}/{\rm d.o.f.} \rangle_{z}$ is the average of $\chi^{2}/{\rm d.o.f.}$ from fitting for each $z$.
}
\label{fig:chisqmap_hyp}
\end{figure*}


\begin{figure*}[tbp]
\centering
\includegraphics[width=8.5cm]{images/Comp_overlap_quark.pdf}
\includegraphics[width=8.5cm]{images/Comp_clover_quark.pdf}
\includegraphics[width=8.5cm]{images/Comp_overlap_Pion.pdf}
\includegraphics[width=8.5cm]{images/Comp_clover_Pion.pdf}
\includegraphics[width=8.5cm]{images/Comp_clover_Nucleon.pdf}
\caption{Comparison between renormalized matrix elements Exp[$g(z)-m_{0}z$] for unsmearing (hyp0, red) and HYP smearing (hyp1, blue) cases.
}
\label{fig:Re_hyp}
\end{figure*}

We start with a fit function with two parameters in the linearly divergent part, $k$ and $\Lambda_{\rm QCD}$. 
Fig.~\ref{fig:chisqmap_hyp} shows the $\chi^{2}$ map with respect to them. It is clear
from the plots that in the ``small-$\chi^2$ band'', we cannot find a solution for $k$ with the  one-loop perturbative value 7.4 GeV$^{-1}$fm$^{-1}$ (1.46 if dimensionless, see Eq.~(\ref{eq:k})) any more. 
This is expected because the lattice space is enlarged by a factor of two, and therefore the values of $k$ in the ``small-$\chi^2$ band'' are expected to be at most
half of this, named 3.7 GeV$^{-1}$fm$^{-1}$ (0.73 if dimensionless). In fact, most probable $k$ is now around 
3.0 GeV$^{-1}$fm$^{-1}$ (0.59 if dimensionless) or below. Another way of saying this is HYP smearing smooths out the short-range behaviors such as the linear divergence. The value of $\Lambda_{\rm QCD}$, on the other hand, has a larger range,
possibly corresponding again to a larger effective lattice spacing.

So for the HYP smearing cases, we choose a set of ($k$, $\Lambda_{\rm QCD}$) near the center of the ``small-$\chi^{2}$ band''. We choose $k=2.5$ GeV$^{-1}$fm$^{-1}$ (0.49 if dimensionless) for overlap quark and clover quark, and $k=2.6$ GeV$^{-1}$fm$^{-1}$ (0.51 if dimensionless) for overlap pion, clover pion and clover nucleon. The corresponding $\Lambda_{\rm QCD}$ shows a larger dispersion, 
with $\Lambda_{\rm QCD} = 0.43$ and $0.38$ GeV for overlap quark and pion correlations respectively, and $\Lambda_{\rm QCD} = 0.61$, $0.37$ and $0.35$ GeV for clover 
quark, pion and nucleon correlations, respectively.

Once the ($k$, $\Lambda_{\rm QCD}$) parameters are chosen, we 
subtract the linear divergences, and match the result
to the perturbative matrix elements to determine possible
further subtractions for the non-perturbative mass
effect. Fig.~\ref{fig:Re_hyp} shows that the renormalized matrix elements for HYP smearing in blue ticks and lines, 
in comparison with the cases without HYP smearing. 

The large difference between smeared and unsmeared matrix
elements is seen in the case of clover quark correlation 
at large $z$. 
Again, we speculate that due to the chiral symmetry breaking effects, the linear divergence has some curious behavior there. However, the difference is smaller in the pion matrix element where the error bars are also bigger. For the overlap case, 
the correlations for quark and pion show no substantial change 
due to smearing effects. On the other hand, the matrix
element in the nucleon case shows a much smaller 
error bar in the smeared case, which demonstrates the power 
of smearing. 

To conclude, our analysis method to renormalize
the linear divergent matrix elements can successfully be used
for smeared matrix elements as well. Smearing does not seem to change the behavior of renormalization qualitatively, but only quantitatively. 
The intrinsic physics does not seem to change under (moderate) smearing,
consistent with the expectation that smearing is an alternative method to reduce the statistical noise in lattice calculations.
